\chapter{Interfaces and data contracts}
\label{sec:contracts}

Chapter~\ref{sec:architecture} argued that the whole proposal is two
interfaces, two adapters, and a database. This chapter makes those
five things exact enough to code against. The discipline throughout:
every field and every method either names the verdict it carries,
with the chapter holding the evidence, or is labeled an engineering
judgment that the build team owns. A reader should be able to point
from any line of the eventual code review back to a line of this
survey, or know that no such line exists.

One vocabulary note before the details. A \emph{study} is the unit a
researcher opens: one search space, one objective declaration, one
budget, many trials. The contracts below are organized the way a
study meets them: first what the researcher declares, then the two
decision-makers (searcher and scheduler), then the store that
remembers everything, then the executors that pay the bills.

\section{The study declaration}

A study begins with a declaration, and the declaration is where
several chapters' verdicts become unavoidable rather than optional.

The \emph{objective declaration} names the quantity being optimized,
its direction, and its evaluation protocol. For RL studies that means
choosing among final evaluated return, area under the training curve,
and return at a fixed budget, explicitly, because
Chapter~\ref{sec:workloads} showed the three orderings genuinely
disagree; the field is mandatory and recorded, never defaulted
silently. Multi-objective studies declare a vector of such
objectives, and the product treats the front as the study's output
from that moment (Chapter~\ref{sec:moo}).

The \emph{seed protocol} is part of the declaration, with defaults
rather than blanks: a number of tuning seeds per candidate and a
disjoint set of test seeds for the final report, the protocol of
\citet{eimer2023hyperparameters} that Chapters~\ref{sec:problem}
and~\ref{sec:workloads} adopted. Opting out is possible and is
itself recorded, so a study that skipped the protocol says so in its
own audit trail.

The \emph{budget declaration} states the study's size in full-run
equivalents and the number of concurrently available full-length
workers. The second number exists because of
Chapter~\ref{sec:population}'s boundary: at eight or more dedicated
workers with a schedule-shaped question the population line becomes
eligible, and below it the product steers to the rung methods. The
declaration is what lets the product enforce that boundary in
software rather than trusting each user's optimism.

\section{The search-space schema}

The schema is a tree of knob declarations, one per hyperparameter,
and its fields track exactly the structure
Chapter~\ref{sec:problem} taught.

\begin{description}[style=nextline, leftmargin=1.6em, itemsep=2pt]
\item[kind] One of continuous, ordinal, categorical. Three values,
not the four kinds of Chapter~\ref{sec:problem}, because the fourth
is structural rather than typed: a conditional knob is a continuous,
ordinal, or categorical knob that also carries a predicate, so
conditionality lives in \textbf{condition} below. Integer-valued
knobs are ordinals, the ordinal kind being defined as an integer with
a natural order, and there is no separate integer kind to choose. The
distinction that remains
is load-bearing downstream: ordinals are modeled as ordered, never as
unordered labels, because respecting the order produced measurable
gains in the BG-PBT ablations (Chapter~\ref{sec:bayesian},
Section~\ref{sec:highdim}), and categorical axes route to overlap
kernels or per-category machinery rather than being embedded and
hoped about.
\item[range and transform] Bounds for continuous and ordinal knobs,
the value list for categorical ones, and an optional transform, with
log-uniform as the expected choice for scale-type knobs; the 90
percent of a uniform draw wasted in the top decade
(Chapter~\ref{sec:modelfree}) is the argument, and the same flag
drives input warping in the GP searcher.
\item[condition] An optional parent-knob-and-value predicate. A
conditional knob exists only when its predicate holds, which makes
the space a tree rather than a box; TPE consumes this natively, and
searchers that cannot must say so through their capability flags
below rather than silently sampling meaningless coordinates
(Chapters~\ref{sec:problem} and~\ref{sec:bayesian}).
\item[prior] An optional distribution over the knob, the expert's
belief about where the optimum lies. First-class from day one
because retrofitting it is expensive and every project lead has one
(Chapter~\ref{sec:transfer}); consumed by the $\pi$BO weight and the
PriorBand portfolio once tier~1 lands, recorded but inert before
that.
\item[note] Free text for units and meaning. An engineering judgment
with a purpose: the schema doubles as the study's documentation, and
Chapter~\ref{sec:problem}'s knob tour is the register this field is
meant to echo.
\end{description}

\section{The searcher interface}

A searcher proposes configurations and learns from results. The
interface is deliberately small.

\begin{description}[style=nextline, leftmargin=1.6em, itemsep=2pt]
\item[propose(n, state)] Return $n$ configurations to evaluate next.
Batch proposals are the normal case, not the extension, because idle
workers were Chapter~\ref{sec:bayesian}'s standing argument; a GP
searcher implements $n > 1$ by joint batch acquisition and handles
still-running trials by fantasies, exactly the machinery of
Section~\ref{sec:acquisition}.
\item[observe(record)] Ingest one trial record from the store,
including partial-fidelity observations; a searcher that only
understands final scores declares that limitation.
\item[capabilities] A static declaration: which knob kinds, whether
conditionals are supported, whether multiple objectives are, whether
a prior is consumed, maximum sensible dimension. The product refuses
a mismatched study loudly at declaration time. The flag exists
because the audit of Chapter~\ref{sec:libraries} found exactly such
mismatches shipping silently (a continuous-only PB2 is the standing
example), and refusing early is cheaper than diagnosing a silently
degenerate study.
\item[snapshot and restore] Serialize internal state, so long
studies survive restarts and population methods can fork searcher
state along a lineage. Engineering judgment on the mechanism; the
requirement itself follows from the population line's design
(Chapter~\ref{sec:population}).
\end{description}

Behind this interface, tier~0 ships wrapped implementations (Optuna
TPE, the BoTorch GP with qLogEI and dimension-scaled priors, CMA-ES,
Sobol and log-uniform random), tier~1 adds qLogNEHVI, Chebyshev
scalarization, the $\pi$BO weight, and EI-per-cost as decorators on
the GP searcher, and the trust-region mixed-space searcher arrives as
the shared core of high-dimensional single-objective search and the
population line's explore step (Chapters~\ref{sec:bayesian}
and~\ref{sec:population}).

\section{The scheduler interface}

A scheduler allocates budget across running trials. Its contract has
one more verb than the classical statement, because the population
line needs it.

\begin{description}[style=nextline, leftmargin=1.6em, itemsep=2pt]
\item[report(trial, fidelity, value)] Called as metrics stream in;
returns continue, stop, or pause. ASHA and the median rule live
entirely behind this verb, which is why they ask nothing of training
code beyond streaming a metric (Chapter~\ref{sec:multifidelity}).
\item[promote()] Return paused or new trials to run when a worker
frees, with rung bookkeeping ($\eta$, initial budget $r$) internal
to the scheduler. Asynchronous promotion against results-so-far is
the ASHA rule; a PASHA flag stops growing the top rung once
rankings stabilize.
\item[gate(predicate)] Register a veto evaluated at every rung. A
trial failing its gate is discarded regardless of its score; this is
the sampler workload's divergence and split-$\widehat{R}$ machinery
generalized into a hook (Chapter~\ref{sec:workloads}), and MO-ASHA's
correctness gates use the same mechanism (Chapter~\ref{sec:moo}).
\item[exploit(loser, winner) and explore(config)] The population
verbs: copy weights and hyperparameters from winner to loser through
the executor's checkpoint surgery, then perturb or re-propose via a
searcher. Keeping these on the scheduler interface, with the
searcher consulted for explore, is precisely the entanglement
Chapter~\ref{sec:population} described, made explicit instead of
implicit.
\end{description}

Every scheduler decision lands in the store with its reason, because
the audit trail is a product feature, not a log file: the researcher
of Chapter~\ref{sec:product}'s walk reads exactly this record.

\section{The run store}

The store is one database with three kinds of rows, and several
chapters' promises are really promises about these rows.

Trial rows carry the configuration, seed, rung history with
per-rung scores and costs, the full learning curve, final objective
values as a vector, veto outcomes, and a parent-trial pointer. The
parent pointer is the population line's lineage: with it, a
hyperparameter \emph{schedule} is a walk up the tree, queryable
after the study, which is what ``schedules are a first-class
output'' meant concretely in Chapter~\ref{sec:population}.

Study rows carry the declaration in full (space, objective, seeds,
budget, boundary inputs), the incumbent or front over time, and the
study-level diagnostics, of which one is computed always: the rank
correlation between rung scores and final scores.
Chapter~\ref{sec:multifidelity} promised that the crossing-curves
hazard would become a measured property of our workloads; this row
is where the measurement lives, and the validation suite consumes
it.

Query surfaces are part of the contract because two verdicts depend
on them: warm starting a new study from configurations that ranked
well on similar past studies (Chapter~\ref{sec:transfer}) is a
ranked query over trial rows, and hypervolume-over-budget reporting
(Chapter~\ref{sec:roadmap}) is a scan over study rows. Storage
technology is an engineering judgment; the schema above is not.

\section{The executor adapters}

Two adapters, per the architecture chapter. The Ray Tune adapter
maps trials to Tune trials and must expose checkpoint surgery as
first-class operations: save, copy weights and optimizer state
across trials, resume. That is the machinery PBT exercises inside
Tune today and the single strongest reason
Chapter~\ref{sec:architecture} kept Tune underneath; our adapter's
job is to keep it reachable from the scheduler's exploit verb. The
local adapter runs minutes-scale trials in-process for the
Hamiltonian regime, where Chapter~\ref{sec:workloads} showed the
distributed stack's overhead is the visible cost. Both adapters
meter per-trial cost themselves, because EI-per-cost consumes
measured cost, not the user's folklore (Chapter~\ref{sec:transfer}).

\section{The architecture factory}
\label{sec:archfactory}

One contract exists only because Chapter~\ref{sec:workloads} decided
that architecture search means width, depth, and a few structural
flags rather than topology. Those coordinates sit in the ordinary
search space, so the searcher and scheduler interfaces above need no
new verbs for them; what they do need is a way to turn a
configuration into a model, and a way to say when two configurations'
models can share weights. Both belong to the user's workload, not to
the tuner, which is why they are a contract rather than an
implementation.

\begin{description}[style=nextline, leftmargin=1.6em, itemsep=2pt]
\item[build(config)] Return an initialized model from the
architecture coordinates of a configuration. The workload supplies
this; the product never constructs architectures itself, because the
space of things a user might mean by ``depth'' is not ours to
enumerate. The factory is also what makes the excluded methods
excluded by construction: a contract that maps coordinates to a model
has no place to put a supernet.
\item[measure(model)] Return the size objectives, parameter count and
FLOPs per forward pass, for a built model. These are measured on the
constructed model rather than predicted from the coordinates, which
is the whole reason Chapter~\ref{sec:workloads} was willing to treat
size as an objective and unwilling to treat proxies as a substitute
for training. Multi-objective studies consume these through the
ordinary objective vector (Chapter~\ref{sec:moo}).
\item[compatible(config$_a$, config$_b$)] Declare whether weights
from one configuration's model can be loaded into another's. Pure
hyperparameter changes are compatible; a width or depth change is
not. The scheduler's exploit verb consults this before checkpoint
surgery, and an incompatible answer routes the trial to the transfer
policy below instead of copying weights into a shape that cannot hold
them. A workload that gets this predicate wrong produces silent
garbage, so the conformance suite tests it directly: every declared
compatible pair must survive a load, and every incompatible pair must
refuse one.
\item[transfer(parent, child)] The policy for an incompatible
inheritance: restart the child from initialization, or warm-start it
by distillation from the parent, which is BG-PBT's route
(Chapter~\ref{sec:population}). The choice is declared per study, not
inferred, and its cost is charged to the trial that caused it rather
than hidden in the population's overhead. Distillation arrives with
the population line in tier~2; before then the only legal policy is
restart, and the factory says so through its capability flags.
\end{description}

The searcher's \textbf{capabilities} declaration gains one flag to
match: whether the searcher models architecture coordinates as
ordinals and categoricals with the rest of the space, or requires
them held fixed. A searcher answering the latter is not eligible for
a study whose space includes architecture axes, and the product
refuses the combination at declaration time for the same reason it
refuses a continuous-only searcher on a categorical space.

\section{What would overturn this chapter}

Contracts are falsified by construction work, and the suite is
arranged to do it early: if the tier~0 wraps cannot be expressed
behind these interfaces without holes, if the reference-parity tests
of Chapter~\ref{sec:validation} need interface changes to pass, or
if the population verbs prove insufficient for a faithful PBT run
on the Ray adapter, the contracts change and this chapter is
rewritten before the code doubles down. The interfaces are the
product's most expensive thing to change later; that is why they
get a chapter, a register mapping, and the first weeks of the
calendar rather than an afternoon.
